\chapter{Cómo usar este libro}

\section*{¿Para quién es?}

Para cualquier persona que quiera entender, desde cero, cómo funciona un dron que vuela solo:
estudiantes de licenciatura, ingenieros de otras áreas o curiosos. No
necesitas saber de drones, de control automático ni de inteligencia artificial. Sí ayuda saber
usar una computadora con Linux y recordar algo de física y matemáticas de bachillerato; cuando
haga falta algo más, lo explicamos en el momento.

El libro cuenta una historia real: cómo se construyó, en tres días de octubre de 2026, una
simulación para el monitoreo de incendios forestales. En ella un dron Tarot Ironman 650
despega, recorre un bosque de 420 por 420 metros, encuentra un incendio con su cámara, avisa
sus coordenadas GPS y se queda vigilándolo. Cada pieza de esa historia es una excusa para explicar
un concepto.

\section*{Cómo está organizado}

\begin{description}
  \item[Parte I. Lo básico.] Qué es un dron y cómo vuela, la física mínima que hace falta, qué es
        una simulación y qué es ROS~2, el ``sistema nervioso'' que conecta todas las piezas.
  \item[Parte II. El control.] Cómo se consigue que el dron vaya exactamente a donde queremos aunque
        sople el viento: control PID, modos deslizantes, super-twisting y el tema central del
        proyecto, la discretización implícita.
  \item[Parte III. Visión.] Qué es una imagen para una computadora, cómo funciona una cámara y cómo
        una red neuronal (YOLO) aprende a reconocer fuego y humo.
  \item[Parte IV. Construir la simulación.] La instalación y cómo se metieron en Gazebo el dron, sus
        motores, el bosque, las montañas, los animales y el incendio.
  \item[Parte V. Autonomía.] Cómo el dron decide qué hacer: recorrer el bosque, evitar árboles y
        montañas, estimar dónde está el fuego, confirmarlo y avisar. También cómo volar con un
        piloto automático real (ArduPilot).
  \item[Parte VI. Manos a la obra.] Un tutorial para hacer tu primer vuelo, la lista de órdenes y
        cómo se verificó que todo funciona.
\end{description}

Al final hay un \textbf{glosario} con todas las palabras raras.

Si solo quieres usar la simulación, puedes ir directo al capítulo~\ref{cap:primer-vuelo}
(``Tu primer vuelo''). Si quieres entenderla, léelo en orden: cada parte se apoya en la anterior.

\section*{Los recuadros}

\begin{recuerda}
Lo más importante del tema, en pocas palabras. Si solo lees esto, ya te llevas la idea.
\end{recuerda}

\begin{consejo}
Trucos prácticos que ahorran tiempo o problemas.
\end{consejo}

\begin{ejemplo}
Una comparación con algo cotidiano para entender un concepto abstracto.
\end{ejemplo}

\begin{tecnico}
Detalles matemáticos o de programación. Si es tu primera lectura, sáltalos sin culpa.
\end{tecnico}

\begin{aviso}
Errores que es fácil cometer y que pueden romper algo o hacerte perder horas.
\end{aviso}

\begin{leccion}
Problemas reales que aparecieron durante el desarrollo y cómo se resolvieron. Son de lo más
valioso del libro: así se trabaja de verdad.
\end{leccion}

\begin{ejercicio}
Algo que puedes hacer en la simulación para comprobar lo que acabas de leer.
\end{ejercicio}

\section*{Convenciones de escritura}

\begin{itemize}
  \item Las órdenes que se escriben en la terminal aparecen en recuadros grises:
\begin{lstlisting}[style=bash]
ros2 launch cuadrirotor_sim simulacion.launch.py
\end{lstlisting}
  \item Los nombres de archivos, programas y tópicos van en letra de máquina: \cod{busqueda_incendio}.
  \item \cod{~} significa tu carpeta personal (\cod{/home/tu_usuario}).
  \item Las distancias están en metros, los ángulos en radianes o grados (se indica) y el tiempo
        en segundos.
\end{itemize}
